% ============================================================
% main.tex
% VVUQ-05 draft-bill (LaTeX)
%   H. R. 9510 BILL v3.0 - a VISUAL DRAFT amendment to the Federal Food, Drug,
%   and Cosmetic Act (21 U.S.C. 301 et seq.), grounded on the Title 21 device
%   provisions current through Public Law 119-93.
%
% WHAT THIS FILE IS. This is the v3.0 visual scaffold for H. R. 9510. It is NOT
% a finished bill. It keeps the full congressional amendment structure of Bill
% v2.0 (the VVUQ-04 final-bill: caption, A BILL, enacting clause, SECTION 1 short
% title and table of contents, SEC. 2 findings, SEC. 3 the operative amendment
% and new section 515D, SEC. 4 comparative print) so it stands on its own as a
% bill, and it ADDS a more visual perspective: every operative part carries a
% figure or table slot and a set-off DRAFTING INSTRUCTIONS block that tells a
% future Claude Code Opus 4.8 (1M context) Max pass exactly which files in this
% repository to process to render that figure or table on a white background as
% a text-based diagram, with no raster image. Five non-operative appendices
% extend the perspective: Appendix A visualizes the VVUQ-01 and VVUQ-02
% engineering evidence; Appendix B instructs the build of every required U.S.
% Bill submission deliverable; Appendix C establishes the VVUQ-01 through VVUQ-04
% developments as a new standard for explainability of the legislation-creation
% process; Appendix D is the research-influence matrix; and Appendix E is the
% development-transparency statement and commit schedule.
%
% WHY A VISUAL PERSPECTIVE. Bill v3.0 answers three goals. First, it is a
% visualization of Bill v2.0, built on the VVUQ-01 through VVUQ-04 record and the
% verbatim visual catalog in papers/VVUQ-05/update-bill/figures-bill. Second, it
% scaffolds the supplementary Markdown deliverables that papers/VVUQ-05/update-
% bill/next-steps says a U.S. House bill needs before submission, to be added in
% a later step under papers/VVUQ-05/deliverables. Third, it establishes the
% four-work evidence-to-law pipeline as a reusable explainability standard. The
% bill is separately compiled in LaTeX into a PDF; the deliverables are separate
% Markdown documents, each complete and submission quality at this bill version.
%
% Compile recipe (Overleaf, pdfLaTeX):
%   pdflatex main.tex
%   bibtex   main
%   pdflatex main.tex
%   pdflatex main.tex
%
% Formatting (centralized in usctitle.sty): white background throughout;
% RaggedRight body with even interword spacing and no right-margin overflow;
% maximal widow and orphan penalties so no single line is stranded; an
% abbreviated running header with no overlapping fields; URL breaking in the
% references; single hyphens only; the section symbol for every codified
% reference; full-width left-aligned ragged-right tables; and ASCII figures set
% verbatim in white, black-ruled monospace frames that preserve spacing exactly.
% ============================================================
\documentclass[11pt]{article}
\usepackage{usctitle}

\title{Verification Before Generation in Physical AI Oncology Trials Act of 2026
(H. R. 9510), Bill v3.0: A Visual Draft Amendment to the Federal Food, Drug, and
Cosmetic Act}
\author{Prepared for VVUQ-05 by CEO Kevin Kawchak, ChemicalQDevice}
\date{June 3, 2026 (current through Public Law 119-93)}

\begin{document}

% ============================================================
% BILL CAPTION AND FRONT MATTER (the caption page). The House and Senate
% statements are condensed: the long title appears once, the referral block is
% brief, and the standard enacting clause is retained. An ASCII overview of the
% v3.0 visual draft process closes the page.
% ============================================================
\thispagestyle{plain}
\pagenumbering{roman}

\billcaption{119TH CONGRESS}{2d Session}{H.\,R. 9510}

\billlongtitle{To amend the Federal Food, Drug, and Cosmetic Act to require an
automated verification, validation, and uncertainty quantification process to
clear robot-patient interaction code before that code is generated or executed
in a Physical AI oncology clinical investigation.}

\billintro{IN THE HOUSE OF REPRESENTATIVES}{June 3, 2026}{%
Mr. \underline{Kevin Kawchak} authored the following bill with the assistance of
artificial intelligence; which was referred to the Committee on Energy and
Commerce}

\billhead

\enactingclause

\billcovernote{Bill v3.0 is a visual draft of H. R. 9510, the Verification
Before Generation in Physical AI Oncology Trials Act of 2026. It keeps the full
operative text of Bill v2.0 so it stands on its own, and it adds a visual
perspective: figure and table slots, an engineering-evidence appendix, a
submission-deliverables appendix, and an explainability-standard appendix. It
was prepared with the assistance of an autonomous artificial intelligence coding
agent (Claude Code Opus 4.8, 1M context, Max) under human direction. It is an
independent research draft, is not an enacted law, and is not legal advice; the
development record appears in Appendix E.}

\begin{asciifig}
  EVIDENCE (VVUQ-01, VVUQ-02)         LAW (VVUQ-03 bill, VVUQ-04 amendment)
  +----------------------------+      +-----------------------------------+
  | method + 51-test pipeline  |      | H. R. 9510 standalone bill        |
  | 10-gate humanoid, 172 test | ===> | FD&C Act amendment, new 360e-5    |
  +----------------------------+      +-----------------------------------+
               |                                       |
               +------------------+--------------------+
                                  v
             VVUQ-05/update-bill  (the inputs to this draft)
             +---------------------------------------------+
             | figures-bill : every ASCII diagram + table  |
             | next-steps   : the submission deliverables  |
             +---------------------------------------------+
                                  |
                                  v
         H. R. 9510 BILL v3.0   <== THIS VISUAL DRAFT SCAFFOLD
         bracketed instructions name the exact files to process
                                  |
                                  v
         A future Claude Code pass renders the figures and the
         deliverables, yielding the final visual bill and the
         papers/VVUQ-05/deliverables submission package
\end{asciifig}
\figcaption{Cover Figure. The Bill v3.0 visual draft process: VVUQ-01 through
VVUQ-04 evidence and law, reorganized through the VVUQ-05 update-bill inputs into
this scaffold, then rendered by a future pass into the final bill and the
deliverables package.}

\clearpage
\pagenumbering{arabic}

% ============================================================
% SECTION 1. SHORT TITLE; TABLE OF CONTENTS.
% ============================================================
\billsec{SECTION 1.}{SHORT TITLE; TABLE OF CONTENTS.}

\uscprov{1}{\textbf{(a) Short Title.} This Act may be cited as the
``Verification Before Generation in Physical AI Oncology Trials Act of 2026''.}

\uscprov{1}{\textbf{(b) Table of Contents.} The table of contents for this Act is
as follows:}

\begin{center}
\begin{tabular}{@{}>{\raggedright\arraybackslash}p{0.9\linewidth}@{}}
Sec. 1. Short title; table of contents.\\
Sec. 2. Findings; the evidence to law record.\\
Sec. 3. Amendment to the Federal Food, Drug, and Cosmetic Act.\\
Sec. 4. Comparative print; changes in existing law.\\
Appendix A. Visual engineering evidence (non-operative).\\
Appendix B. Required submission deliverables (non-operative).\\
Appendix C. The VVUQ explainability standard (non-operative).\\
Appendix D. Research influence matrix (non-operative).\\
Appendix E. Development transparency and commit schedule.\\
\end{tabular}
\end{center}

\begin{draftbox}{Section 1 - short title and table of contents}
\dstep{Retain the short title verbatim from papers/VVUQ-04/final-bill/main.tex (SECTION 1(a)); verify the bill name and the Act name against papers/VVUQ-04/instruct-bill/10-legal-crosswalk-and-bill-style.md (Section D and the drafting-conventions checklist in Section E, item 1: verify every bill and law name). Confirm H. R. 9510 and "119TH CONGRESS, 2d Session" against papers/VVUQ-04/final-bill/README.md.}
\dstep{Keep the table of contents as a single-column left-aligned list (the ragcol helper at width 0.9 of the line); list Sec. 1 through Sec. 4 and Appendix A through Appendix E exactly as above; keep each entry to one line and reword any wrap so no entry leaves a one- or two-word last line.}
\dstep{Do not number the new FD&C Act section here; the new section number (FD&C Act section 515D, 21 U.S.C. section 360e-5) is established in Sec. 3 and recorded in the clerical amendment to the Act's own table of contents.}
\end{draftbox}

% ============================================================
% BILL BODY AND APPENDICES (one file per section; each corresponds to main.tex).
% The bill body (SEC. 2 through SEC. 4) carries the operative text and the
% comparative print; the appendices carry the visual, deliverable, and
% explainability perspective. Each file opens with its own heading and ends with
% the DRAFTING INSTRUCTIONS that name the exact files a future pass must process.
% ============================================================
%% ============================================================
%% sections/s2-findings.tex
%% H. R. 9510 Bill v3.0 - SEC. 2. FINDINGS; THE EVIDENCE TO LAW RECORD.
%% Common visuals in this section: Figure 1 (four-work evidence-to-law lineage),
%% Table 1 (the four works at a glance), Figure 2 (accelerated timeline). The
%% findings are presented as a condensed narrative plus the three visuals, so the
%% section uses prose and figures in place of a long numbered list.
%% ============================================================
\billsec{SEC. 2.}{FINDINGS; THE EVIDENCE TO LAW RECORD.}

\uscprov{1}{\textbf{(a) In General.} Congress finds that an automated process
must verify robot-patient interaction code before that code is generated or
executed in a Physical AI oncology clinical investigation, and that the record of
the VVUQ-01 through VVUQ-04 developments, visualized in this section, establishes
both the need for and the feasibility of that requirement \cite{kawchak2026national,
vvuq05-figures-bill}.}

\figslot{Figure 1.}{Four-Work Evidence-to-Law Lineage}{VVUQ-01 method and
pipeline, to VVUQ-02 hard engineering proof, to VVUQ-03 proposed bill, to VVUQ-04
precise amendment; evidence flows forward and each work cites the records of the
prior works.}

\uscprov{1}{\textbf{(b) The Embodiment Problem and the Sequencing Gap.} Physical
AI oncology investigations place autonomous and semi-autonomous robots in direct
physical contact with patients, and the behavior of a robot at the moment of
contact is determined by software that is increasingly machine-generated. In a
disembodied pipeline an error yields a wrong number a human can catch downstream;
in an embodied surgical system the same error class yields a wrong motion a human
cannot catch in time. No Federal law requires the verification of robot-patient
interaction code to clear before that code is generated or executed, because the
order of operations is not regulated \cite{kawchak2026vvuq01, kawchak2026national}.}

\figslot{Table 1.}{The Four Works at a Glance}{Each work, what it is, its release
range, its headline assurance result (51 of 51 tests; 172 of 172 tests and
composite 93.56; 15 sections and 21 tables; new section 360e-5 and 10 conforming
amendments), and its end-goal DOI.}

\uscprov{1}{\textbf{(c) Feasibility Is Demonstrated.} Automated verification of
robot-patient interaction code is feasible today. A pipeline held the
verification fraction at 1.0, accepting one candidate artifact and blocking five,
and a ten-gate assurance suite over control software for a surgical humanoid
cleared reproducibly from a fixed seed, each gate bound to a published external
standard \cite{kawchak2026vvuq01, kawchak2026vvuq02}.}

\uscprov{1}{\textbf{(d) The Closest Existing Analogues.} The predetermined change
control plan authority in section 515C of the Federal Food, Drug, and Cosmetic
Act (21 U.S.C. 360e-4) and the Food and Drug Administration final guidance on such
plans for artificial-intelligence-enabled device software functions are the
closest existing analogues to validating a change before deployment, but each
remains voluntary, submission-based, and neither automated nor required in
advance of generation. The Quality Management System Regulation (part 820 of
title 21, Code of Federal Regulations), effective February 2, 2026, and the
electronic-records requirements of part 11 of such title supply quality and
record-integrity scaffolding but do not impose a rule that verification precede
generation \cite{fda-pccp-2024, fr-qmsr-2024}.}

\uscprov{1}{\textbf{(e) A Recognized Standards Basis.} Recognized consensus
standards, including ASME V\&V 40-2018, the Food and Drug Administration guidance
assessing the credibility of computational modeling, NASA-STD-7009A, IEEE Std
7009-2024, IEC 80601-2-77, ISO 14971, and ISO/TS 15066, supply the credibility
framework, dispersion bounds, fail-safe requirements, and biomechanical contact
limits to which an automated gate may be bound \cite{asmevv40, fda-credibility-2023,
nasastd7009, isots15066}.}

\uscprov{1}{\textbf{(f) Human-Over-AI Review Is Already Widely Adopted.}
Human-over-AI review is already widely adopted in State law, including the
California Physicians Make Decisions Act (Senate Bill 1120), the Illinois Wellness
and Oversight for Psychological Resources Act (House Bill 1806), the Texas
Responsible Artificial Intelligence Governance Act (House Bill 149), and Maryland
House Bill 820, and the strongest Federal precedent for mandatory human-over-AI
oversight is the Medicare Advantage artificial-intelligence guardrail (section
422.101(c) of title 42, Code of Federal Regulations) \cite{ca-sb1120, il-hb1806,
tx-hb149, md-hb820, cms-4201-ma}.}

\figslot{Figure 2.}{Accelerated Timeline Versus Conventional Methods}{The
conventional months-to-years path for a method paper, verified robotics code, a
manuscript, and a federal bill, set beside the about eleven days of autonomous,
real-time work that produced VVUQ-01 through VVUQ-04.}

\uscprov{1}{\textbf{(g) Modernization, Nondiscrimination, and Classification.} As
clinical investigation moves toward real-time and decentralized designs governed
by harmonized good clinical practice, manual review cannot keep pace with the
volume of machine-generated control code, and an automated gate that clears
before generation serves that modernization. Federal nondiscrimination law,
including section 92.210 of title 45, Code of Federal Regulations, title VI of the
Civil Rights Act of 1964 (42 U.S.C. 2000d), and the Americans with Disabilities
Act of 1990 (42 U.S.C. 12101 et seq.), requires that a patient-facing algorithm
minimize the risk of bias and respect accessibility, and the artificial-
intelligence taxonomy maintained by the American Medical Association provides a
graduated spine for matching verification rigor to the degree of autonomy of a
device \cite{iche6r3, cfr-92-210, usc-titlevi, usc-ada, cpt-appendix-s}.}

\begin{draftbox}{Section 2 - findings and the three findings visuals}
\dstep{Figure 1 (four-work evidence-to-law lineage): render the lineage diagram verbatim from papers/VVUQ-05/update-bill/figures-bill/01RootREADME.md (Visual 3, under the heading "Building process from one work to the next"), reconciled against the same diagram in papers/README.md and the "Lineage" diagram in papers/VVUQ-04/final-bill/README.md. Set it inside the asciifig environment so the boxes, arrows, and DOIs keep their exact spacing on a white background; place it as a full-width text figure. Do not redraw it as a raster image.}
\dstep{Table 1 (the four works at a glance): build the table from papers/VVUQ-05/update-bill/figures-bill/01RootREADME.md (Visual 2) and papers/README.md ("The four works at a glance"). Use the xltabular Y and L columns at the body measure with columns Work, What it is, Releases, Headline assurance result, and End-goal DOI; keep every cell left aligned and ragged right (each width prepended with \rabs) so no big interword gaps open; reword long cells onto two lines rather than letting any cell overflow.}
\dstep{Figure 2 (accelerated timeline): render the conventional-versus-this side-by-side timeline verbatim from papers/VVUQ-05/update-bill/figures-bill/01RootREADME.md (Visual 5) and papers/README.md ("Accelerated timeline versus conventional methods") in the asciifig environment; confirm the per-work spans (VVUQ-01 about 5 days, VVUQ-02 about 2 days, VVUQ-03 about 1 day, VVUQ-04 about 2 days, about eleven days total).}
\dstep{Findings prose: carry the fourteen grounded findings of Bill v2.0 from papers/VVUQ-04/final-bill/main.tex (SEC. 2, paragraphs (1) through (14)) and condense them into the lettered narrative paragraphs above, preserving every cited authority but using prose and the three visuals in place of a long numbered list, consistent with the instruction to use fewer numbered and lettered lists. Keep each finding tied to a recorded authority and resolve every citation to a key in references.bib.}
\dstep{Grounding: re-confirm the technical findings against papers/VVUQ-04/instruct-bill/09-vvuq-standards-clinical-trial-oncology.md and papers/VVUQ-03/final-paper/sections/definitions.tex; the legal-context findings against papers/VVUQ-04/instruct-bill/01-federal-statutory-framework.md, 02-fda-ai-device-regulation.md, 04-cms-coverage-payment-ai.md, and 05-privacy-security-nondiscrimination.md; and the state human-in-the-loop practice against 06-state-medical-ai-laws.md. Keep the 119th Congress emerging bills and the executive actions out of the findings and confine them to Appendix D per 10-legal-crosswalk-and-bill-style.md (Section B).}
\dstep{Formatting: single hyphens only; the section symbol for every codified reference; even RaggedRight spacing; reword so no finding ends in a one- or two-word line and no single line is stranded onto the next page; keep the white background on both figures and the table.}
\end{draftbox}

%% ============================================================
%% sections/s3-amendment.tex
%% H. R. 9510 Bill v3.0 - SEC. 3. AMENDMENT TO THE FEDERAL FOOD, DRUG, AND
%% COSMETIC ACT. Carries the operative new section 515D (21 U.S.C. 360e-5), the
%% ten conforming amendments, the clerical amendment, the rule of construction,
%% and the effective date, so the bill stands on its own.
%% Common visuals: Table 2 (ten-gate threshold schedule), Figure 3 (gate
%% decision surface and funnel), Figure 4 (statutory layering through Title 21),
%% Table 3 (ten conforming amendments crosswalk).
%% ============================================================
\billsec{SEC. 3.}{AMENDMENT TO THE FEDERAL FOOD, DRUG, AND COSMETIC ACT.}

\uscprov{1}{\textbf{(a) In General.} Subchapter V of the Federal Food, Drug, and
Cosmetic Act (21 U.S.C. 351 et seq.) is amended by inserting after section 515C
(21 U.S.C. 360e-4) the following new section:}

\uscquote{0}{``SEC. 515D. VERIFICATION BEFORE GENERATION OF ROBOT-PATIENT
INTERACTION CODE IN PHYSICAL AI ONCOLOGY INVESTIGATIONS.}
\uscquote{0}{``(a) Requirement.- No robot-patient interaction code may be
generated or executed in a Physical AI oncology clinical investigation unless an
automated verification, validation, and uncertainty quantification process, bound
to one or more named external standards, has first cleared for that interaction,
and the cleared verification record has been documented and attested under
subsection (e).}
\uscquote{0}{``(b) Order of Operations.- A sponsor shall not generate, and shall
not execute, robot-patient interaction code in advance of the cleared
verification record. A record created after generation or execution does not
satisfy this section.}
\uscquote{0}{``(c) Gate Thresholds and Decision Rule.- For each robot-patient
interaction, the verification process shall apply the gate thresholds set out in
the schedule in Table 2. Each gate shall emit a determination of ACCEPT, BLOCK,
or ESCALATE; a single failing dimension shall set BLOCK, and ambiguity or
divergence shall set ESCALATE and return the interaction to a qualified human.
For each interaction the verification fraction shall equal 1.0, the validation
agreement shall meet or exceed the per-gate threshold with the maximum relative
error at or below the per-gate bound and the human review recorded, the
coefficient of variation across not fewer than three seeded runs shall not exceed
the per-gate bound, and, for an interaction subject to a hard catastrophe
predicate, that predicate shall also hold \cite{asmevv40, nasastd7009}.}

\figslot{Table 2.}{Ten-Gate Threshold Schedule}{For each gate: the minimum
validation agreement, the maximum coefficient of variation, whether the gate
carries a hard catastrophe predicate, and the governing external standards.}

\figslot{Figure 3.}{VVUQ Gate Decision Surface and Funnel}{The
ACCEPT, BLOCK, ESCALATE decision rule over the verify, validate, quantify
dimensions, set beside the candidate-to-ACCEPT funnel that accepted one artifact
and blocked five.}

\uscquote{0}{``(d) Readiness Gates.- Before a patient is enrolled or a robot
acts, the trial site shall pass the Physical AI Standard Level gate on
legislative authorization, regulatory compliance, and operational readiness; each
robot shall meet the minimum Unification Standard Level for its functional type,
which is not less than 7.0 for a surgical type, 6.0 for a therapeutic positioning
or diagnostic needle placement type, 4.0 for a rehabilitative type, and 3.0 for a
companion monitoring type; and each robot's behavior shall be validated across not
fewer than two simulation frameworks within the cross-framework tolerances of
less than 2 millimeters of trajectory deviation and less than 0.5 newton of force
discrepancy \cite{kawchak2026national}.}
\uscquote{0}{``(e) Documentation and Attestation.- Before robot-patient
interaction code is generated or executed, the sponsor shall make and file, for
each covered algorithm, algorithm documentation including a plain-language
summary, the decision logic or architecture, the training or conditioning data
sets, the gate bindings identifying the governing external standard and the
threshold enforced, the verification-before-generation record, and a determinism
statement giving the seed; a compliance statement and a signed attestation by a
responsible officer; and a hash-chained, tamper-evident audit trail that satisfies
part 11 of title 21, Code of Federal Regulations \cite{cfr-part11, cfr-hipaa-164}.}
\uscquote{0}{``(f) Cybersecurity, Human Oversight, and Lifecycle.- Each Physical
AI system shall incorporate authentication and access control, encryption of data
in transit and at rest, network segmentation, and software-integrity
verification; shall operate under recorded human oversight with a manual override
and an emergency stop that brings the system to a safe state within the
procedure-appropriate budget and a hand-back-to-human escalation when a gate
escalates or a fault is detected; and shall maintain configuration and
model-version management, monitor for model and data drift, and, on
decommissioning, archive the audit trail, securely erase patient data, and revoke
access \cite{iec60601, iso13849}.}
\uscquote{0}{``(g) Nondiscrimination.- Each covered algorithm and its training
data shall minimize the risk of bias on protected characteristics and adhere to
evidence-based clinical guidelines, and each Physical AI system shall be designed
and operated consistent with applicable accessibility and accommodation
requirements \cite{cfr-92-210, usc-titlevi, usc-ada}.}
\uscquote{0}{``(h) Regulations.- Not later than 365 days after the date of
enactment of this section, the Secretary shall issue final regulations to carry
out this section. Pending those regulations, the Secretary may issue interim
guidance, and a sponsor may comply by meeting the requirements of subsections (a)
through (g).}
\uscquote{0}{``(i) Definitions.- In this section, the terms `Physical AI system',
`robot agent', `robot-patient interaction code', `generative artificial
intelligence', `AI-enabled device software function', `verification',
`validation', `uncertainty quantification', `verification fraction', `validation
agreement', `coefficient-of-variation bound', `hard catastrophe predicate',
`Physical AI Standard Level', `Unification Standard Level', and `hash-chained
audit trail' have the meanings carried forward from the prior bill and reconciled
with section 201(h) and section 520(o) of this Act.}
\uscquote{0}{``(j) Rule of Construction.- Nothing in this section displaces the
investigational-device exemption under section 520(g), the protection of human
subjects under part 50 of title 21, Code of Federal Regulations, the
investigational new drug framework under part 312 of such title, or the
predetermined change control plan authority under section 515C as it applies
outside Physical AI, and nothing in this section itself reclassifies any
device.''.}

\uscprov{1}{\textbf{(b) Conforming Amendments.} The Federal Food, Drug, and
Cosmetic Act is amended, with the corresponding changes shown in the comparative
print in section 4, so that the device definition in section 201(h) reaches
software that generates or executes robot-patient interaction code; the prohibited
acts in section 301 reach a violation of section 515D; a device used without a
cleared section 515D record is adulterated under section 501; the real world
evidence program in section 505F recognizes section 515D records; a
verification-governed change consistent with a predetermined change control plan
needs no new premarket notification under section 510(k); special controls under
section 513 are graduated by degree of autonomy; a premarket approval application
under section 515 must include the cleared verification record; a predetermined
change control plan under section 515C must include an automated
verification-before-change protocol; the clinical-decision-support exclusion in
section 520(o) does not reach robot-control software; and a savings clause at
section 521 preserves State human-in-the-loop requirements that are not different
from, or in addition to, a Federal device requirement \cite{usc-321h, usc-331,
usc-351, usc-355g, usc-360, usc-360c, usc-360e, usc-360e4, usc-360j, usc-360k}.}

\figslot{Figure 4.}{Statutory Layering of New Section 360e-5 Through Title 21}{The
device-regulation pathway the amendment threads: is it a device (321(h)); not
excluded by the CDS carve-out (360j(o)); classify (360c); 510(k) or PMA (360,
360e); the predetermined change control plan keystone (360e-4); and the new,
mandatory, ex ante section 360e-5.}

\figslot{Table 3.}{Ten Conforming Amendments Crosswalk}{Each conforming amendment:
the section file, the 21 U.S.C. citation, the FD\&C Act section, and the change made
by H. R. 9510.}

\uscprov{1}{\textbf{(c) Clerical Amendment.} The table of contents for the Federal
Food, Drug, and Cosmetic Act is amended by inserting after the item relating to
section 515C the following new item:}
\uscquote{0}{``Sec. 515D. Verification before generation of robot-patient
interaction code in Physical AI oncology investigations.''.}

\uscprov{1}{\textbf{(d) Rule of Construction.} Nothing in this Act displaces the
investigational-device exemption under section 520(g) of the Federal Food, Drug,
and Cosmetic Act (21 U.S.C. 360j(g)), the protection of human subjects under part
50 of title 21, Code of Federal Regulations, the investigational new drug
framework under part 312 of such title, the predetermined change control plan
authority under section 515C as it applies outside Physical AI, or any State law
preserved under section 521 as amended; and nothing in this Act reclassifies any
device \cite{cfr-part50, cfr-part312}.}

\uscprov{1}{\textbf{(e) Effective Date; Implementation Deadline; Transition Rule.}
This Act takes effect on the first day of the first calendar quarter beginning not
less than one year after enactment; the Secretary shall issue the final
regulations required by section 515D(h) not later than 365 days after enactment;
and a Physical AI oncology clinical investigation already enrolling on the
effective date shall have 180 days to come into compliance with section 515D, and
this Act does not disturb the February 2, 2026, effective date of the Quality
Management System Regulation \cite{fr-qmsr-2024}.}

\begin{draftbox}{Section 3 - the operative amendment and its four visuals}
\dstep{Table 2 (ten-gate threshold schedule): render the ten-gate table from papers/VVUQ-04/final-bill/main.tex (the xltabular under section 515D(c)(4): gate, minimum agreement, maximum CV, hard predicate, governing external standards), reconciled against papers/VVUQ-05/update-bill/figures-bill/03VVUQ02.md (the "Ten VVUQ Gate Funnel" and the per-gate primary-standards and master tables) and papers/VVUQ-03/final-paper/sections/statutory\_text.tex. Use the L and Y columns (each width prepended with \raggedright\arraybackslash) at the body measure; reword long standard lists onto two lines so no cell overflows; keep the three hard-predicate gates (06 vascular no-fly, 09 human collision, 10 fault and emergency stop) at agreement 1.00.}
\dstep{Figure 3 (gate decision surface and funnel): render two synthesized text panels side by side in the asciifig environment. The decision surface (verify, validate, quantify dimensions to the ACCEPT, BLOCK, ESCALATE rule) comes from papers/VVUQ-05/update-bill/figures-bill/02VVUQ01.md ("The three dimensions, as executed" and the gate decision-surface table) and 03VVUQ02.md ("The three dimensions, as executed"). The funnel (candidates to ACCEPT, with the 6-to-1 VVUQ-01 case and the 10-gate VVUQ-02 case) comes from 02VVUQ01.md (the VVUQ gate funnel, 6 candidates to 1 accepted) and 03VVUQ02.md (the Ten VVUQ Gate Funnel). Preserve original spacing on a white background.}
\dstep{Figure 4 (statutory layering): render the device-pathway diagram verbatim from papers/VVUQ-05/update-bill/figures-bill/05VVUQ04.md (the "How the sections interact" and "How the conforming amendments thread through Title 21" diagrams) and the same diagram in papers/VVUQ-04/final-bill/README.md; mark the new section 360e-5 node and the section 360e-4 keystone clearly. Set it in the asciifig environment at full width.}
\dstep{Table 3 (ten conforming amendments crosswalk): build from papers/VVUQ-04/final-bill/README.md ("Per-section comparative print (changes in existing law)") with columns File, 21 U.S.C., FD\&C section, and Change made by H. R. 9510, for s301, s321, s331, s351, s355g, s360, s360c, s360e, s360e-4, s360j, and s360k; left aligned and ragged right at the body measure.}
\dstep{Operative text: carry the new section 515D and the ten conforming amendments verbatim in substance from papers/VVUQ-04/final-bill/main.tex (SEC. 3); keep the (a) through (j) subsection lettering because statutory form requires it, but render the long subsection (i) definitions and the conforming list as condensed prose rather than expanded numbered lists. Anchor each duty to an in-force authority per papers/VVUQ-04/instruct-bill/10-legal-crosswalk-and-bill-style.md (Section B): the change-control link to section 360e-4 and the FDA PCCP final guidance (keys usc-360e4, fda-pccp-2024); the credibility basis to ASME V\&V 40-2018 (key asmevv40); record integrity to 21 CFR part 11 and the HIPAA Security Rule (keys cfr-part11, cfr-hipaa-164); nondiscrimination to 45 CFR 92.210, Title VI, and the ADA (keys cfr-92-210, usc-titlevi, usc-ada); and human oversight to the CMS Medicare Advantage guardrail (key cms-4201-ma).}
\dstep{Formatting: single hyphens only; the section symbol for every codified reference (correct any "SS" to the section symbol); even RaggedRight spacing with no overflow; place each figure or table slot near the subsection it supports and avoid leaving a single subsection line stranded on the next page; keep all four visuals on a white background.}
\end{draftbox}

%% ============================================================
%% sections/s4-comparative.tex
%% H. R. 9510 Bill v3.0 - SEC. 4. COMPARATIVE PRINT; CHANGES IN EXISTING LAW.
%% Presents the comparative print visually: the convention statement, a
%% per-section change map (Table 4), and the device-pathway threading figure
%% (Figure 5). The eleven marked section files are reproduced by the future pass
%% from the Bill v2.0 sources named in the draftbox.
%% ============================================================
\billsec{SEC. 4.}{COMPARATIVE PRINT; CHANGES IN EXISTING LAW.}

\uscprov{1}{The following comparative print shows the changes proposed to existing
law by this Act. Matter proposed to be deleted is shown as \cprstrike{matter to be
deleted}; matter proposed to be inserted is shown as \cprinsert{matter to be
inserted}; existing law in which no change is proposed is shown as regular text,
and unchanged matter omitted for brevity is shown by a row of asterisks
\cite{usc-fdca, olrc-usc21}.}

\usccrosshead{CHANGES IN EXISTING LAW MADE BY THE BILL}
\begin{center}\small\itshape Federal Food, Drug, and Cosmetic Act (reproduced from
Title 21,\\United States Code, current through Public Law 119-93)\end{center}

\uscprov{1}{The comparative print reaches eleven Title 21 device sections in 21
U.S.C. order, each shown only as to its affected provisions. The map in Table 4
identifies every reproduced section, its 21 U.S.C. citation, the Federal Food,
Drug, and Cosmetic Act section number, and the change H. R. 9510 makes to it, and
Figure 5 shows where each change sits on the device-regulation pathway the new
section 360e-5 attaches to \cite{pl-119-93}.}

\figslot{Table 4.}{Per-Section Comparative-Print Change Map}{For each reproduced
section (s301, s321, s331, s351, s355g, s360, s360c, s360e, s360e-4, s360j,
s360k): the 21 U.S.C. citation, the FD\&C Act section, and the marked change.}

\figslot{Figure 5.}{How the Eleven Reproduced Sections Thread Through Title
21}{The device-regulation pathway with the changed nodes marked: the device
definition and CDS exclusion, classification, the 510(k) and PMA pathways, the
predetermined change control plan keystone, and the cross-cutting adulteration,
prohibited-acts, real-world-evidence, and preemption duties.}

\uscprov{1}{Each reproduced section, with its insertions and deletions marked,
follows in the rendered bill in the order s301, s321, s331, s351, s355g, s360,
s360c, s360e, s360e-4, s360j, and s360k \cite{usc-360e4, usc-360j, usc-360k}.}

\begin{draftbox}{Section 4 - comparative print, the change map, and the threading figure}
\dstep{Table 4 (per-section change map): build from papers/VVUQ-04/final-bill/README.md ("Per-section comparative print (changes in existing law)") with columns File, 21 U.S.C., FD\&C section, and Change made by H. R. 9510, for all eleven sections; use the L and Y columns at the body measure (each width prepended with \rabs); keep cells left aligned with no overflow and no big interword gaps.}
\dstep{Figure 5 (threading figure): render the device-pathway diagram from papers/VVUQ-05/update-bill/figures-bill/05VVUQ04.md ("How the conforming amendments thread through Title 21") and the same diagram in papers/VVUQ-04/final-bill/README.md ("How the conforming amendments thread through Title 21"); this is the same pathway as Figure 4 but annotated to the comparative print, so keep the two visually consistent and mark which nodes carry a marked change. Set it verbatim in the asciifig environment on a white background.}
\dstep{Marked comparative print: reproduce the eleven affected sections, wrapping each insertion in the cprinsert marker and each deletion in the cprstrike marker (the same two markers shown in the convention statement above), from the finished section files in papers/VVUQ-04/final-bill/sections/ (s301.tex, s321.tex, s331.tex, s351.tex, s355g.tex, s360.tex, s360c.tex, s360e.tex, s360e-4.tex, s360j.tex, s360k.tex), which already carry the focused affected-provisions print and the asterisk convention for omitted matter. Wire them in with one LaTeX input line per section in 21 U.S.C. order, exactly as papers/VVUQ-04/final-bill/main.tex does, and keep each change consistent with the operative amendatory language in SEC. 3.}
\dstep{Source of record: confirm the reproduced text against papers/VVUQ-04/template-bill/xml\_usc21@119-93.zip (the OLRC USLM XML) and name Public Law 119-93 only as the currency point, never as the law being amended, per papers/VVUQ-04/instruct-bill/10-legal-crosswalk-and-bill-style.md.}
\dstep{Formatting: single hyphens only; the section symbol for every codified reference; even RaggedRight spacing with no overflow; keep the bracket convention (no strikethrough or underline) and avoid stranding a single reproduced line onto the next page; keep the table and figure on a white background.}
\end{draftbox}

%% ============================================================
%% sections/a5-evidence.tex
%% H. R. 9510 Bill v3.0 - APPENDIX A. VISUAL ENGINEERING EVIDENCE (NON-OPERATIVE).
%% The visualization of Bill v2.0: the VVUQ-01 and VVUQ-02 engineering record that
%% grounds the findings, drawn from papers/VVUQ-05/update-bill/figures-bill.
%% Common visuals: Figure 6 (gate funnel comparison), Figure 7 (procedure
%% timelines), Figure 8 (humanoid safety figure), Table 5 (external standards map).
%% ============================================================
\clearpage
\phantomsection
\addcontentsline{toc}{section}{Appendix A. Visual Engineering Evidence (Non-Operative)}
\usccrosshead{APPENDIX A. VISUAL ENGINEERING EVIDENCE (NON-OPERATIVE)}

\uscprov{1}{This appendix is non-operative. It visualizes the engineering record
of VVUQ-01 and VVUQ-02 that grounds the findings in section 2 and the gate
discipline codified in section 515D, so a reader can see the evidence the bill
rests on without leaving the document. None of these figures is cited as the
legal basis of an operative clause; only the enacted statutes, in-force rules, and
recognized standards named in the bill carry operative weight
\cite{kawchak2026vvuq01, kawchak2026vvuq02}.}

\figslot{Figure 6.}{The Gate, From One Dimension to Ten}{Side by side: the VVUQ-01
six-candidate funnel that accepted one artifact and blocked five over the verify,
validate, quantify dimensions, and the VVUQ-02 ten-gate funnel that cleared a
surgical humanoid behavior against ten standards-bound gates, three of them
carrying a hard catastrophe predicate.}

\uscprov{1}{The gate is the load-bearing object of the whole record: VVUQ-01 holds
it to a single behavior across three dimensions, and VVUQ-02 scales it to ten
humanoid-specific gates without changing the ACCEPT, BLOCK, ESCALATE semantics.
The verification fraction stays an equality at 1.0 in both, which is why the bill
can require it as an equality \cite{kawchak2026vvuq01, kawchak2026vvuq02}.}

\figslot{Figure 7.}{The Procedure Timelines}{The VVUQ-02 sixty-second eight-phase
Whipple performed by one humanoid with two hands, set over the VVUQ-01 hundred-
and-sixty-eight-day PDAC hybrid pilot, so the single-procedure assurance window
and the full adjuvant pilot window are read on one page.}

\figslot{Figure 8.}{The Humanoid Safety Figure}{A combined view of the
seventy-one-degree-of-freedom kinematic chain, the balance zero-moment-point and
support polygon, and the vessel no-fly proximity bands, the three drawn artifacts
that make the catastrophe gates concrete.}

\uscprov{1}{The safety figure is where embodiment becomes visible: the kinematic
chain shows how many ways a body can move, the support polygon shows the balance
margin a fall would cross, and the no-fly bands show the vascular volume the hard
predicate forbids breaching. These are the dimensions a wrong motion would violate
and a human could not catch in time \cite{isots15066, iec8060127}.}

\figslot{Table 5.}{External Standards Map}{The recognized consensus standards bound
to the gates: ASME V\&V 40, NASA-STD-7009A, IEEE 7009, IEC 80601-2-77, ISO 14971,
ISO/TS 15066, ISO 13482, IEC 60601-1, ISO 13849-1, and ISO 9283, with the gate or
duty each governs.}

\begin{draftbox}{Appendix A - the four visual-evidence artifacts}
\dstep{Figure 6 (gate funnel comparison): render two funnels side by side in the asciifig environment. The VVUQ-01 funnel (6 candidates to 1 accepted, with the per-stage verification, validation, and uncertainty cuts) comes verbatim from papers/VVUQ-05/update-bill/figures-bill/02VVUQ01.md (the "VVUQ Gate Decision Funnel" image spec and the gate decision-surface table). The VVUQ-02 ten-gate funnel (candidates to ACCEPT, with the per-gate V and CV thresholds and the three immediate-catastrophe gates) comes verbatim from papers/VVUQ-05/update-bill/figures-bill/03VVUQ02.md (the "Ten VVUQ Gate Funnel (candidates to ACCEPT)" block). Preserve original spacing on a white background.}
\dstep{Figure 7 (procedure timelines): render the sixty-second eight-phase Whipple timeline from papers/VVUQ-05/update-bill/figures-bill/03VVUQ02.md ("60-Second 8-Phase Whipple Timeline") over the one-hundred-sixty-eight-day PDAC hybrid pilot timeline from papers/VVUQ-05/update-bill/figures-bill/02VVUQ01.md (the PDAC pilot timeline and pdac\_hybrid\_pilot\_timeline.txt block). Keep both as fenced ASCII timelines in the asciifig environment; do not redraw as a chart.}
\dstep{Figure 8 (humanoid safety figure): combine the three drawn artifacts from papers/VVUQ-05/update-bill/figures-bill/03VVUQ02.md - the "H2-Surgical 1.0 Kinematic Chain (71 DOF)", the "Balance ZMP and Support Polygon", and the "Vessel No-Fly Proximity Bands (VVUQ 06)" - into one full-page text figure, each panel labeled, all in the asciifig environment with exact spacing on a white background.}
\dstep{Table 5 (external standards map): build from papers/VVUQ-05/update-bill/figures-bill/03VVUQ02.md (the "Per-gate primary standards", "Cross-cutting standards", and "standards-map" tables) reconciled with the ten-gate schedule in papers/VVUQ-04/final-bill/main.tex (section 515D(c)(4)); columns Standard, What it governs, and Gate or duty; left aligned and ragged right at the body measure (each width prepended with \raggedright\arraybackslash).}
\dstep{Synthesis guidance: these four artifacts are the visual basis for the findings in section 2 and the gate schedule in section 515D, so keep the numbers consistent across them (verification fraction 1.0; the three hard-predicate gates 06, 09, 10; composite mean 93.56; seed 20260525). Use the correlation map in papers/VVUQ-05/update-bill/figures-bill/README.md (Section 4, the cross-file map and the shared evidentiary objects) to confirm each artifact's source and that the 10-gate schedule is the same object cited as evidence in VVUQ-03 and anchored to Title 21 in VVUQ-04.}
\dstep{Formatting: single hyphens only; the section symbol where a codified reference appears; even RaggedRight spacing; keep every figure and table on a white background; prefer full vertical page figures and avoid large empty white bands by sizing each asciifig to its content.}
\end{draftbox}

%% ============================================================
%% sections/a6-deliverables.tex
%% H. R. 9510 Bill v3.0 - APPENDIX B. REQUIRED SUBMISSION DELIVERABLES
%% (NON-OPERATIVE). The instruction set for the supplementary Markdown documents a
%% future pass writes under papers/VVUQ-05/deliverables for U.S. House submission,
%% grounded in papers/VVUQ-05/update-bill/next-steps.
%% Common visual: Table 6 (the submission deliverables matrix).
%% ============================================================
\clearpage
\phantomsection
\addcontentsline{toc}{section}{Appendix B. Required Submission Deliverables (Non-Operative)}
\usccrosshead{APPENDIX B. REQUIRED SUBMISSION DELIVERABLES (NON-OPERATIVE)}

\uscprov{1}{This appendix is non-operative. It identifies the supplementary
documents a draft United States House bill needs before it can be introduced, as
recorded in the pre-submission analysis in papers/VVUQ-05/update-bill/next-steps,
and it instructs a future pass to write each one as a separate, complete, and
submission-quality Markdown document under papers/VVUQ-05/deliverables. The bill
itself is the LaTeX measure compiled separately into a PDF; these deliverables are
the package that accompanies it to a sponsoring office \cite{vvuq05-next-steps,
crs-r44001}.}

\uscprov{1}{The two reality checks from the next-steps analysis govern the whole
package: H. R. 9510 is illustrative because the Clerk assigns a bill its number
only at introduction, and a private project cannot submit a bill directly because
only a sitting Member may introduce a House bill. Every deliverable is therefore
framed to get the text into final form and into the hands of a Member who will
introduce it, and each preserves the standing disclaimer that the work is an
independent research draft, not pending legislation, and not legal advice
\cite{house-rule-xii, congress-howlawsmade}.}

\figslot{Table 6.}{The Submission Deliverables Matrix}{For each deliverable: the
output Markdown file under papers/VVUQ-05/deliverables, its purpose, and the lead
source files a future pass synthesizes it from.}

\begin{draftbox}{Appendix B - Table 6 and the per-deliverable build instructions}
\dstep{Table 6 (the deliverables matrix): build a left-aligned, ragged-right xltabular at the body measure (each width prepended with \rabs) with columns Deliverable and output file, Purpose, and Lead source files, with one row per deliverable below. State at the top that every deliverable is a standalone Markdown document, complete and submission quality at this bill version, with white-background Markdown tables, single hyphens, the section symbol where a codified reference appears, and no images or mermaid diagrams.}
\dstep{Deliverable 01 - papers/VVUQ-05/deliverables/01-one-page-summary.md: a one-page plain-English summary of H. R. 9510 (the problem, the verify-before-generate rule, the new section 515D, the ten conforming amendments, and the assurance numbers). Synthesize from papers/VVUQ-04/final-bill/main.tex (SECTION 1 and SEC. 2), papers/VVUQ-04/final-bill/README.md ("The bill at a glance"), and papers/README.md ("The four works at a glance"). Keep it to one page and submission quality.}
\dstep{Deliverable 02 - papers/VVUQ-05/deliverables/02-section-by-section-analysis.md: a section-by-section analysis of every operative part. Synthesize from papers/VVUQ-04/final-bill/main.tex (SEC. 3, the section 515D subsections (a) through (j), and the ten conforming amendments) and papers/VVUQ-04/final-bill/README.md ("The new operative section" and "Per-section comparative print"). One subheading per subsection and per conforming amendment, each with a plain-English explanation; complete and submission quality.}
\dstep{Deliverable 03 - papers/VVUQ-05/deliverables/03-plain-english-policy-memo.md: a problem-and-solution policy memo for a Member's staff. Synthesize from papers/VVUQ-03/final-paper/sections/policy\_memo.tex and problem\_statement.tex and papers/VVUQ-04/instruct-bill/01-federal-statutory-framework.md through 05-privacy-security-nondiscrimination.md. State the sequencing gap, why now, and the in-force analogues (PCCP, QMSR, the CMS guardrail); complete and submission quality.}
\dstep{Deliverable 04 - papers/VVUQ-05/deliverables/04-legislative-findings.md: the standalone findings of fact. Synthesize from papers/VVUQ-04/final-bill/main.tex (SEC. 2, the fourteen findings) and this bill's sections/s2-findings.tex, each finding tied to its references.bib authority; complete and submission quality.}
\dstep{Deliverable 05 - papers/VVUQ-05/deliverables/05-ramseyer-comparative-print.md: the Ramseyer comparative print (changes to existing law) in Markdown. Synthesize from papers/VVUQ-04/final-bill/sections/ (all eleven section files) and this bill's sections/s4-comparative.tex, showing deletions and insertions for each affected section; complete and submission quality.}
\dstep{Deliverable 06 - papers/VVUQ-05/deliverables/06-constitutional-authority-statement.md: the Constitutional Authority Statement required by House Rule XII, clause 7(c). Cite the Commerce Clause (Article I, section 8, clause 3) on which the Federal Food, Drug, and Cosmetic Act rests; synthesize the form and currency from papers/VVUQ-04/instruct-bill/01-federal-statutory-framework.md and resolve to references.bib keys crs-if12335 and house-rule-xii; complete and submission quality.}
\dstep{Deliverable 07 - papers/VVUQ-05/deliverables/07-paygo-and-cost-estimate.md: a PAYGO and cost-estimate note and an earmark declaration. Frame the budgetary effects of the section 515D(h) rulemaking and the Secretary's implementation, drawing on papers/VVUQ-04/final-bill/main.tex (SEC. 3(e) and section 515D(h)) and references.bib key cbo-paygo; note that formal scoring follows introduction; complete and submission quality.}
\dstep{Deliverable 08 - papers/VVUQ-05/deliverables/08-sponsor-and-cosponsor-packet.md: the sponsor and original cosponsor packet (the one sponsor, the cosponsor form, and the Energy and Commerce referral expectation). Synthesize from papers/VVUQ-05/update-bill/next-steps/output-1-next-steps.md (Phase 2) and output-2-next-steps.md (steps 4 and 7) and references.bib keys crs-rs22477 and house-ehopper; complete and submission quality.}
\dstep{Deliverable 09 - papers/VVUQ-05/deliverables/09-stakeholder-engagement-plan.md: the stakeholder engagement and outreach plan (FDA and HHS, patient-advocacy groups, industry, academic societies, and the standards bodies). Synthesize from papers/VVUQ-05/update-bill/next-steps/output-1-next-steps.md (Phase 2 stakeholder engagement) and papers/VVUQ-04/instruct-bill/09-vvuq-standards-clinical-trial-oncology.md (the standards bodies: ASME V\&V 40, IEC 80601-2-77, ISO 13482, ISO/TS 15066); keep outreach respectful and forward looking; complete and submission quality.}
\dstep{Deliverable 10 - papers/VVUQ-05/deliverables/10-legislative-counsel-routing-memo.md: the House Office of Legislative Counsel routing memo (legislative form, amendatory instructions, cross-references, positive-law question). Synthesize from papers/VVUQ-05/update-bill/next-steps/output-2-next-steps.md (steps 2 and 3) and papers/VVUQ-04/instruct-bill/10-legal-crosswalk-and-bill-style.md; resolve to references.bib key house-holc; complete and submission quality.}
\dstep{Deliverable 11 - papers/VVUQ-05/deliverables/11-currency-and-cross-reference-matrix.md: the currency and cross-reference verification matrix (re-verify that no newer public law has changed Title 21, that the target number section 515D / 21 U.S.C. section 360e-5 is still unused, and that every cross-reference and each of the ten conforming amendments still lands). Synthesize from papers/VVUQ-05/update-bill/next-steps/output-1-next-steps.md (Phase 1 currency check), papers/VVUQ-04/instruct-bill (current through May 31, 2026), and references.bib key pl-119-93; a left-aligned matrix; complete and submission quality.}
\dstep{Deliverable 12 - papers/VVUQ-05/deliverables/12-testimony-and-research-influence-brief.md: hearing testimony and the research-influence brief, the one place the 119th Congress emerging bills and the executive actions may be discussed (never in the bill text). Synthesize from papers/VVUQ-04/instruct-bill/07-executive-actions-national-ai-strategy.md and 08-emerging-federal-bills.md and this bill's sections/a8-research-matrix.tex; complete and submission quality.}
\dstep{Deliverable 00 - papers/VVUQ-05/deliverables/README.md: the package index (badges, a repository-structure tree of the deliverables folder, a table of contents, and a one-line gloss per deliverable), consistent with this draft-bill README; no images or mermaid diagrams; white-background tables; single hyphens.}
\dstep{Cross-cutting requirements: every deliverable is its own Markdown file, each complete and submission quality at this bill version (Bill v3.0), so a sponsoring office could act on it without the others. Keep the bill itself in LaTeX (compiled to PDF separately) and never fold a deliverable into the bill text. Preserve the disclaimers, use single hyphens only, render every table on a white background, use the section symbol for codified references, and include no images, mermaid diagrams, or colored figures.}
\end{draftbox}

%% ============================================================
%% sections/a7-explainability.tex
%% H. R. 9510 Bill v3.0 - APPENDIX C. THE VVUQ EXPLAINABILITY STANDARD
%% (NON-OPERATIVE). Establishes the VVUQ-01 through VVUQ-04 developments as a new
%% standard for explainability of the U.S. legislation-creation process.
%% Common visuals: Figure 9 (prompt evolution across the works), Figure 10
%% (internal bill evolution), Table 7 (the explainability standard).
%% ============================================================
\clearpage
\phantomsection
\addcontentsline{toc}{section}{Appendix C. The VVUQ Explainability Standard (Non-Operative)}
\usccrosshead{APPENDIX C. THE VVUQ EXPLAINABILITY STANDARD (NON-OPERATIVE)}

\uscprov{1}{This appendix is non-operative. It establishes the VVUQ-01 through
VVUQ-04 developments as a reusable standard for explaining how a piece of United
States legislation was created. The claim is narrow and demonstrable: every step
that produced this bill, from the first method paper to the precise amendment,
recorded the exact prompt that drove it, the full output it produced, a VVUQ gate
decision that held the output to its specification, and a commit trail that pushed
each file in real time. A reader can therefore reconstruct how the law was built,
which is the property an explainable process must have \cite{nist-ai-explainability,
nist-ai-600}.}

\figslot{Figure 9.}{Prompt Evolution Across the Works}{The chain of recorded
prompts from VVUQ-02 and VVUQ-03 (instruct, draft, full) through VVUQ-04 (instruct,
update, draft, full) to VVUQ-05 (figures, next-steps, this draft), showing how the
instruction to the agent sharpened as the work moved from method to statute.}

\uscprov{1}{The prompt record is the first half of the standard. Because each
stage keeps its prompt and its output side by side, the reasoning that turned
evidence into statutory language is inspectable rather than hidden: a reviewer can
read prompt-a-update-bill and see the instruction that retargeted the bill onto
the Federal Food, Drug, and Cosmetic Act, then read the output that resulted
\cite{vvuq05-figures-bill}.}

\figslot{Figure 10.}{Internal Bill Evolution}{The VVUQ-04 internal stages
instruct-bill, template-bill, update-bill, draft-bill, full-bill, and final-bill,
extended by the VVUQ-05 update-bill (figures-bill and next-steps) into this v3.0
draft-bill, each stage a labeled box feeding the next.}

\uscprov{1}{The stage record is the second half. The bill did not appear in one
pass: a legal crosswalk was assembled, the real statute was reproduced, the
amendment was scaffolded with bracketed instructions, the instructions were
executed, and the result was finalized, with this visual draft a further stage.
Each stage is a directory a reviewer can open, so the provenance of every clause
is traceable to the files it was synthesized from \cite{kawchak2026vvuq04bill}.}

\figslot{Table 7.}{The VVUQ Explainability Standard}{Each explainability dimension
(explanation, meaningful, accuracy, knowledge limits), what makes the
legislation-creation step satisfy it, and the artifact in the repository that
records it.}

\uscprov{1}{Mapped to the four principles of explainable artificial intelligence,
the standard reads as follows: explanation is met because every step has a
recorded prompt and output; meaningful is met because each stage carries a README
that explains the step to a human; accuracy is met because the VVUQ gate holds
each output to its specification and the results reproduce from a fixed seed; and
the knowledge-limits principle is met because the gate escalates to a human and the
disclaimers mark where the system must hand back. This is the explainability the
legislation-creation process should carry, demonstrated end to end
\cite{nist-ai-explainability, kawchak2026vvuq01}.}

\begin{draftbox}{Appendix C - the explainability standard and its two evolution figures}
\dstep{Figure 9 (prompt evolution): render a left-to-right box-and-arrow chain in the asciifig environment from the recorded prompt files, in this order: papers/VVUQ-02/instructions/prompt-instruct.md, papers/VVUQ-02/draft-paper/prompt-draft-paper.md, papers/VVUQ-02/full-paper/prompt-full-paper.md; papers/VVUQ-03/instruct-paper/prompt-instruct-paper.md, papers/VVUQ-03/draft-paper/prompt-draft-paper.md, papers/VVUQ-03/full-paper/prompt-full-paper.md; papers/VVUQ-04/instruct-bill/prompt-instruct-bill.md, papers/VVUQ-04/update-bill/prompt-update-bill.md and prompt-a-update-bill.md, papers/VVUQ-04/draft-bill/prompt-draft-bill.md, papers/VVUQ-04/full-bill/prompt-full-bill.md; and papers/VVUQ-05/update-bill/figures-bill/prompt-figures-bill.md, papers/VVUQ-05/update-bill/next-steps/prompt-1-next-steps.md and prompt-2-next-steps.md, papers/VVUQ-05/draft-bill/prompt-draft-bill.md. Label each box with the stage and a few words of how the instruction sharpened; keep exact spacing on a white background.}
\dstep{Figure 10 (internal bill evolution): render a box-and-arrow chain in the asciifig environment over the VVUQ-04 stage directories papers/VVUQ-04/instruct-bill, template-bill, update-bill, draft-bill, full-bill, and final-bill, extended by papers/VVUQ-05/update-bill (figures-bill and next-steps) to papers/VVUQ-05/draft-bill (this v3.0 draft). Reuse the lineage motif from papers/VVUQ-05/update-bill/figures-bill/05VVUQ04.md ("Purpose and place in the lineage") and papers/VVUQ-04/final-bill/README.md ("Lineage"); mark the v3.0 endpoint with a "<== THIS DRAFT" pointer.}
\dstep{Table 7 (the explainability standard): build a left-aligned, ragged-right table at the body measure (each width prepended with \rabs) with columns Dimension, What makes the legislation-creation step satisfy it, and Recording artifact. Use the four NISTIR 8312 principles (explanation, meaningful, accuracy, knowledge limits) for the rows; resolve to references.bib keys nist-ai-explainability and nist-ai-600. Map each row to the concrete artifacts: the prompt and output Markdown pairs at each stage; the per-stage README; the VVUQ gate decision surface and the seed-20260525 reproducibility; and the ESCALATE rule and the disclaimers.}
\dstep{Framing source: ground the standard in papers/VVUQ-05/update-bill/figures-bill/README.md (Section 4, how the five catalogs correlate, and Section 5, how a downstream AI should use the corpus) and papers/README.md ("The processing feat accomplished by AI"), which already describe the self-similar, fully recorded pipeline. Keep the claim narrow and demonstrable: this is an explainability standard for the legislation-creation process, not a claim about the merits of the bill.}
\dstep{Formatting: single hyphens only; the section symbol where a codified reference appears; even RaggedRight spacing with no overflow; keep both figures and the table on a white background; size each asciifig to its content so no large empty white band opens, and avoid leaving a single line stranded onto the next page.}
\end{draftbox}

%% ============================================================
%% sections/a8-research-matrix.tex
%% H. R. 9510 Bill v3.0 - APPENDIX D. RESEARCH INFLUENCE MATRIX (NON-OPERATIVE).
%% The 119th Congress emerging bills and the executive actions, kept out of the
%% operative text and confined to this matrix, a memo, testimony, or the findings
%% narrative. Common visual: Table 8 (the research influence matrix).
%% ============================================================
\clearpage
\phantomsection
\addcontentsline{toc}{section}{Appendix D. Research Influence Matrix (Non-Operative)}
\usccrosshead{APPENDIX D. RESEARCH INFLUENCE MATRIX (NON-OPERATIVE)}

\uscprov{1}{The authorities in this appendix are research influences only. They
inform the legislative findings, a policy memorandum, testimony, and this matrix,
but none is cited as the legal basis of an operative clause of this Act. Only
enacted statutes, in-force rules, and recognized standards carry operative weight,
so the 119th Congress emerging bills and the executive actions are kept here and
out of the bill text, consistent with the operative-versus-influence rule the
prior bill follows \cite{kawchak2026vvuq04bill}.}

\figslot{Table 8.}{Research Influence Matrix}{Each research-influence authority:
what it is, what it informs in the bill package, and its placement (memo,
appendix, testimony, or this matrix), with no row cited from an operative clause.}

\uscprov{1}{The executive actions in the matrix include the order on removing
barriers to American leadership in artificial intelligence, the order ensuring a
national policy framework, the Office of Management and Budget memorandum on
accelerating Federal use of artificial intelligence, and the HHS artificial-
intelligence strategy; the 119th Congress bills include the Healthy Technology
Act, the Health Tech Investment Act, the Ban AI Denials in Medicare Act, the
HEALTH AI Act, and the House resolution on the CMS WISeR pilot. Each is named by
number and subject and is verified before use \cite{eo-14179, eo-14365, omb-m2521,
hhs-ai-strategy-2025, hr238-healthy-tech, s1399-health-tech-investment,
hr6361-ban-ai-denials, hr5045-health-ai, hres694-wiser}.}

\begin{draftbox}{Appendix D - the research influence matrix}
\dstep{Table 8 (research influence matrix): build a left-aligned, ragged-right xltabular at the body measure (each width prepended with \raggedright\arraybackslash) with columns Authority, What it is, What it informs, and Placement, carried from papers/VVUQ-04/final-bill/main.tex (Appendix A, the research-influence matrix) and grounded in papers/VVUQ-04/instruct-bill/07-executive-actions-national-ai-strategy.md (EO 14179; EO 14365; OMB M-25-21; HHS AI Strategy) and 08-emerging-federal-bills.md (H.R. 238 Healthy Technology Act; S. 1399 Health Tech Investment Act; H.R. 6361 Ban AI Denials in Medicare Act; H.R. 5045 HEALTH AI Act; H.Res. 694 on the CMS WISeR pilot). State at the top that these are research influences only and never appear in operative text.}
\dstep{Verification caution: process the verification notes in papers/VVUQ-04/instruct-bill/08-emerging-federal-bills.md (for example, confirm there is no VALID Act in the 119th Congress and confirm each sponsor and status before use) and papers/VVUQ-04/instruct-bill/10-legal-crosswalk-and-bill-style.md (Section B, operative law versus research influence; Section E, items 2 to 4). Resolve each entry to a references.bib key in the executive-actions-and-emerging-bills set and do not cite any of these keys from the operative sections 1 through 4.}
\dstep{Placement consistency: keep this matrix consistent with the testimony and research-influence brief in Appendix B (papers/VVUQ-05/deliverables/12-testimony-and-research-influence-brief.md), so the same authority carries the same placement in both; the bill text and the comparative print remain free of these authorities.}
\dstep{Formatting: single hyphens only; the section symbol where a codified reference appears; even RaggedRight spacing with no overflow; keep the table on a white background and reword long cells onto two lines rather than letting any cell overflow.}
\end{draftbox}

%% ============================================================
%% sections/a9-transparency.tex
%% H. R. 9510 Bill v3.0 - APPENDIX E. DEVELOPMENT TRANSPARENCY AND COMMIT SCHEDULE.
%% How Bill v3.0 was built and how a future pass finishes it.
%% Common visuals: Table 9 (version lineage), Table 10 (single-PR commit
%% schedule), Figure 11 (the build-and-handoff process).
%% ============================================================
\clearpage
\phantomsection
\addcontentsline{toc}{section}{Appendix E. Development Transparency and Commit Schedule}
\usccrosshead{APPENDIX E. DEVELOPMENT TRANSPARENCY AND COMMIT SCHEDULE}

\uscprov{1}{\textbf{What was built, and how.} Bill v3.0 (v3.0.0) is the visual
draft of H. R. 9510. It was produced by an autonomous artificial intelligence
coding agent, Claude Code Opus 4.8 (1M context, Max), operating under human
direction in a managed, ephemeral cloud container. It was authored across
sequential commits in a single pull request, one file per commit pushed to GitHub
in real time, with a second-to-last error-fix-and-bundle commit and a final
repository-updates commit. Only kevinkawchak/cancer-automated was edited
\cite{anthropic-claude-code, repo-cancer-automated}.}

\figslot{Table 9.}{Bill Version Lineage}{Bill v1.0 (VVUQ-03 standalone bill), Bill
v2.0 (VVUQ-04 FD\&C Act amendment, v2.0.0 through v2.3.0), and Bill v3.0 (this
VVUQ-05 visual draft, v3.0.0), each with its directory, role, and DOI.}

\uscprov{1}{\textbf{What this version adds.} From Bill v2.0, this version keeps
the full operative text and reorganizes it as a visual perspective: it adds the
figure and table slots and their build instructions, the visual engineering
evidence in Appendix A, the submission deliverables in Appendix B, the
explainability standard in Appendix C, and the commit record here, while confining
the emerging bills and executive actions to Appendix D, exactly as the prior bill
does \cite{kawchak2026vvuq04bill}.}

\figslot{Table 10.}{Single-Pull-Request Commit Schedule}{The ordered commits that
build this bundle: main.tex, the style file, the bibliography, the eight section
files, the README, the error-fix-and-bundle commit with the prompt, output, and
zip, and the final repository-updates commit.}

\uscprov{1}{\textbf{How to finish this measure.} A future Claude Code Opus 4.8 (1M
context) Max pass executes the bracketed drafting instructions in each section
against the named files, rendering every figure and table on a white background as
a text-based diagram and writing each Appendix B deliverable as a separate,
submission-quality Markdown document under papers/VVUQ-05/deliverables. The handoff
is deliberate: this draft names the exact files so the next pass spends its effort
on synthesis, not on searching \cite{vvuq05-figures-bill, vvuq05-next-steps}.}

\figslot{Figure 11.}{The Build-and-Handoff Process}{This v3.0 draft, with its
bracketed instructions and figure and table slots, handed to a future pass that
renders the figures and writes the deliverables, yielding the final visual bill and
the papers/VVUQ-05/deliverables package.}

\begin{draftbox}{Appendix E - the lineage, the commit schedule, and the handoff figure}
\dstep{Table 9 (bill version lineage): build a left-aligned, ragged-right table at the body measure (each width prepended with \rabs) with columns Version, Directory, Role, and DOI. Rows: Bill v1.0, papers/VVUQ-03/final-paper, standalone VVUQ Physical AI Oncology Trial Bill, 10.5281/zenodo.20454870; Bill v2.0, papers/VVUQ-04/final-bill, FD\&C Act amendment (new section 360e-5 and ten conforming amendments), 10.5281/zenodo.20485580; Bill v3.0, papers/VVUQ-05/draft-bill, this visual draft, DOI to be minted. Confirm the directories and DOIs against papers/README.md ("Final Papers and Bills on Zenodo") and papers/VVUQ-04/final-bill/README.md.}
\dstep{Table 10 (commit schedule): build a left-aligned, ragged-right table at the body measure with columns Commit, File or action, and Role. List the fourteen commits of this single pull request in order: (1) main.tex; (2) usctitle.sty; (3) references.bib; (4) sections/s2-findings.tex; (5) sections/s3-amendment.tex; (6) sections/s4-comparative.tex; (7) sections/a5-evidence.tex; (8) sections/a6-deliverables.tex; (9) sections/a7-explainability.tex; (10) sections/a8-research-matrix.tex; (11) sections/a9-transparency.tex; (12) README.md; (13) the error-fix-and-bundle commit adding prompt-draft-bill.md, output-draft-bill.md, and draft-bill-LaTeX.zip; and (14) the final repository-updates commit. Confirm against the actual commit history of this pull request.}
\dstep{Figure 11 (build-and-handoff): render a compact box-and-arrow process in the asciifig environment showing this draft (instructions plus slots) handed to a future pass that renders the figures and writes the deliverables, yielding the final bill and the deliverables package. Reuse the shape of the cover figure in papers/VVUQ-05/draft-bill/main.tex so the two are consistent; keep exact spacing on a white background.}
\dstep{Implementation note: summarize how the measure is implemented as law from papers/VVUQ-04/final-bill/main.tex (Appendix B, the implementation pathway) so the transparency statement and the implementation guidance stay aligned: final regulations within 365 days under section 515D(h), an effective date one year out, a 180-day transition for enrolling investigations, and the readiness gates before enrollment.}
\dstep{Formatting: single hyphens only; the section symbol where a codified reference appears; even RaggedRight spacing with no overflow; keep both tables and the figure on a white background; avoid leaving a single line stranded onto the next page and avoid large empty white bands.}
\end{draftbox}


% ============================================================
% AUTHORSHIP, NOTICES, AND PROVENANCE.
% ============================================================
\clearpage
\phantomsection
\addcontentsline{toc}{section}{Authorship, Notices, and Provenance}
\usccrosshead{AUTHORSHIP, NOTICES, AND PROVENANCE}

\uscnotepar{Prepared by CEO Kevin Kawchak
(\href{https://orcid.org/0009-0007-5457-8667}{ORCID 0009-0007-5457-8667}),
ChemicalQDevice (\href{mailto:kevink@chemicalqdevice.com}{kevink@chemicalqdevice.com}).
Adapted using Claude Code Opus 4.8 (1M context) Max; research assisted by ChatGPT
and Google Gemini.}

\uscnotepar{Notice: This measure was adapted from documents of the United States
Government in the public domain under 17 U.S.C. 105. The reproduced statutory
text is in the public domain; the authoritative version is the United States Code
as published by the Office of the Law Revision Counsel. This work is not endorsed
or sponsored by the FDA, HHS, the OLRC, CFR, ICH, or any member of Congress.}

\uscnotepar{Disclaimer: This work is an independent research draft, is not an
enacted law, and is not legal advice; it is not endorsed, sponsored, or approved
by any trial sponsor, contract research organization, site, institutional review
board, regulator, or medical society. All supporting numbers are simulation or
illustrative results reproducible from the deterministic seed 20260525, and no
patient-identifiable information is included.}

\uscnotepar{Rights: The reproduced statutory text is in the public domain; the
generated amendment framing, figures, tables, and notices are released under the
Creative Commons Attribution 4.0 International License (CC BY 4.0),
\href{https://creativecommons.org/licenses/by/4.0/}{https://creativecommons.org/licenses/by/4.0/}.}

\href{https://doi.org/10.5281/zenodo.20485580}{DOI prior bill (v2.0): 10.5281/zenodo.20485580}

\href{https://doi.org/10.5281/zenodo.20372501}{DOI repository: 10.5281/zenodo.20372501}

\uscnotepar{San Diego. June 3, 2026.}

% ============================================================
% PROVENANCE AND RESEARCH SOURCES (BIBLIOGRAPHY).
% A real congressional bill carries no bibliography; this research-aid back
% matter does, so the operative text stays clean while sources remain auditable.
% ============================================================
\clearpage
\phantomsection
\addcontentsline{toc}{section}{Provenance and Research Sources}
\usccrosshead{Provenance and Research Sources}
\begingroup
\RaggedRight
\nocite{*}
\bibliographystyle{ieeetr}
\bibliography{references}
\endgroup

\end{document}
